\documentclass[a4paper,5pt]{article}

% ======================
% Encoding and language
% ======================
\usepackage[utf8]{inputenc}
\usepackage[T1]{fontenc}
\usepackage[french]{babel}

% ======================
% Layout and graphics
% ======================
\usepackage{geometry}
\usepackage{graphicx}
\usepackage{titlesec}
\usepackage{fancyhdr}
\usepackage{hyperref}
\usepackage{caption}
\usepackage{amsmath}
\usepackage{float}
\usepackage{geometry}
\usepackage{graphicx}
\usepackage{titlesec}
\usepackage{fancyhdr}
\usepackage{hyperref}
\usepackage{caption}
\usepackage{amsmath}
\usepackage{float}
\usepackage{booktabs}
\usepackage{tikz}
\usetikzlibrary{arrows.meta}

% ======================
% Page geometry
% ======================
\geometry{
    a4paper,
    top=25mm,
    bottom=20mm,
    textwidth=185mm
}

\setlength{\parindent}{5pt}
\setlength{\parskip}{0pt}

% ======================
% Section formatting
% ======================
\titleformat{\section}
{\normalfont\bfseries}
{\thesection.}{0.5em}{}

\titleformat{\subsection}
{\normalfont\itshape}
{\thesubsection.}{0.5em}{}

% ======================
% Header / footer
% ======================
\pagestyle{fancy}
\fancyhf{}

\rhead{\includegraphics[scale=0.4]{images/hepia_simple.png}}
\rfoot{\textit{\today}}
\cfoot{\thepage}

\renewcommand{\headrulewidth}{0pt}
\renewcommand{\footrulewidth}{0pt}

% ======================
% Caption style
% ======================
\captionsetup{
    font=footnotesize,
    labelsep=colon
}

% ======================
% Document
% ======================
\begin{document}

\begin{center}
\rule{\linewidth}{0.4pt}
{\large
Rapport de laboratoire\\
\vspace{2mm}
Dimensionnement d'un régulateur pour un vérin hydraulique
}

\vspace{8mm}

Guilhem Rozier Vilardell\footnote{\textit{\href{mailto:guilhem.roziervi@hes-so.ch}{MT2, guilhem.roziervi@hes-so.ch}}},
Noah Michelot\footnote{\textit{\href{mailto:noah.michelot@hes-so.ch}{MT2, noah.michelot@hes-so.ch}}}

\vspace{3mm}

\footnotesize{
\textit{Haute école de paysagerie, d'ingénierie et d'architecture, Genève, Suisse}
}
\rule{\linewidth}{0.4pt}
\vspace{9mm}

\end{center}

\setlength{\parskip}{10pt}


% =================================================
\section{Introduction}

Ce laboratoire a pour objet le dimensionnement et l'implémentation d'un
régulateur PI pour la commande en position d'un plateau mobile entraîné par un
vérin hydraulique, lui-même alimenté par une servovalve. Le régulateur est
dimensionné par la méthode du \emph{approximation des asymptotes} (point~A), puis vérifié en simulation (point~C). Il est ensuite implanté sur l'installation réelle, pilotée par un microcontrôleur ESP32 programmé depuis un schéma \textit{Simulink}, avec une période d'échantillonnage $T_E = 1~\mathrm{ms}$ (point~D). Les observations faites sur l'installation,
en particulier la saturation de la commande, permettent enfin de discuter
l'emploi d'un régulateur PID ou P (point~E).

% =================================================
\section{Modèle du système}

Le modèle de l'installation fourni dans l'énoncé est représenté à la
figure~\ref{fig:boucle}. Il comprend la servovalve $G_{sva}$, le vérin
$G_{vey}$, un gain de fuite $K_{per}$ qui modélise une perturbation de débit,
ainsi que la mesure de position $G_{mes}$~:

\begin{equation}
G_{vey}(s) = \frac{8{,}162 \times 10^{3}}{s}\ [\mathrm{m/(m^3/s)}]
\quad,\quad
G_{sva}(s) = \frac{6{,}05 \times 10^{-6}}{1 + s\,T_2}\ [\mathrm{(m^3/s)/V}]
\quad,\quad
G_{mes}(s) = \frac{50}{1 + s\,T_1}\ [\mathrm{V/m}]
\end{equation}

avec $T_2 = 3~\mathrm{ms}$ et $T_1 = 5~\mathrm{ms}$. Dans l'installation réelle
comme en simulation, la consigne de position (en mètres) est convertie en
volts par le gain $K_{mes} = 50~\mathrm{V/m}$, puis traverse un filtre de
consigne $G_{lc}(s)$ dont le rôle est expliqué au point~A.

\begin{figure}[H]
\centering
\resizebox{\linewidth}{!}{%
\begin{tikzpicture}[>=Stealth, font=\small,
  blk/.style={draw, minimum height=8mm, minimum width=12mm, inner sep=3pt},
  sum/.style={draw, circle, inner sep=0pt, minimum size=5mm}]
  \node (r)    at (0,0)    {$r$~[m]};
  \node[blk] (kmes) at (1.9,0)  {$K_{mes}$};
  \node[blk] (glc)  at (4.0,0)  {$G_{lc}$};
  \node[sum] (s1)   at (6.8,0)  {};
  \node[blk] (gr)   at (8.7,0)  {$G_R$};
  \node[blk] (gsva) at (11.6,0) {$G_{sva}$};
  \node[sum] (s2)   at (14.8,0) {};
  \node[blk] (gvey) at (16.9,0) {$G_{vey}$};
  \node (y)    at (19.3,0) {$y$~[m]};
  \node[blk] (gmes) at (11.6,-2) {$G_{mes}$};
  \node (v)    at (11.6,2)  {$v$~[V]};
  \node[blk, dashed] (kper) at (14.8,2) {$K_{per}$};

  \draw[->] (r)--(kmes);
  \draw[->] (kmes)--(glc);
  \draw[->] (glc)--node[above]{$w$~[V]}(s1);
  \draw[->] (s1)--(gr);
  \draw[->] (gr)--node[above]{$u_{cm}$~[V]}(gsva);
  \draw[->] (gsva)--node[above]{$u$~[m$^3$/s]}(s2);
  \draw[->] (s2)--(gvey);
  \draw[->] (gvey)--(y);
  \fill (18.0,0) circle (1.2pt);
  \draw[->] (18.0,0) |- (gmes);
  \draw[->] (gmes) -| node[pos=0.22, above]{$y_m$~[V]}(s1);
  \draw[->, dashed] (v)--(kper);
  \draw[->, dashed] (kper)-- node[left]{$v$~[m$^3$/s]}(s2);

  \node[font=\scriptsize] at (6.25,0.32)  {$+$};
  \node[font=\scriptsize] at (7.1,-0.6)   {$-$};
  \node[font=\scriptsize] at (14.25,0.32) {$+$};
  \node[font=\scriptsize] at (15.1,0.6)   {$-$};
\end{tikzpicture}}
\caption{Boucle de régulation. $G_R$ désigne le régulateur et $G_{lc}$ le filtre
de consigne. La perturbation de débit $K_{per}$ (en pointillés) n'est pas
simulée au point~C.}
\label{fig:boucle}
\end{figure}

Le système à régler (de $w$ à $y_m$, sans le régulateur) s'écrit~:

\begin{equation}
G_s(s) = G_{sva}(s) \cdot G_{vey}(s) \cdot G_{mes}(s)
= \frac{K_s}{s\,(1+s\,T_2)(1+s\,T_1)}
= \frac{2{,}47}{s\,(1+s\,0{,}003)(1+s\,0{,}005)}
\end{equation}

avec $K_s = 6{,}05 \times 10^{-6} \times 8{,}162 \times 10^{3} \times 50 \approx 2{,}47$.
Ce système comporte un \textbf{intégrateur} (provenant de $G_{vey}$) ainsi que
deux petites constantes de temps parasites, $T_1$ et $T_2$. Il s'agit donc d'un
système de \textbf{type~1}.

% =================================================
\section{Vérin hydraulique}

% -------------------------------------------------
\subsection{Question A: Régulateur PI dimensionné par le réglage en correspondance}

\textbf{Structure du PI.} Le régulateur PI est mis sous la forme~:
\begin{equation}
G_R(s) = \frac{1 + s\,T_n}{s\,T_i} = \frac{1}{s\,T_i} + \frac{T_n}{T_i}
\qquad\Rightarrow\qquad
\text{partie intégrale:} \ \ K_i = \frac{1}{T_i}
\qquad
\text{partie proportionnelle:} \ \ K_p = \frac{T_n}{T_i}
\end{equation}

\textbf{Choix du modèle du système.} Le système à régler $G_s$ contient déjà un
intégrateur et nous souhaitons utiliser un régulateur PI. D'après la méthode
vue en cours, le système doit alors se mettre sous la forme
$G_s(s) = K_s\,/\,\big[s\,(1+s\,T_p)\big]$. Il faut donc regrouper les deux
petites constantes de temps en une constante équivalente $T_p$, en posant
$(1+sT_1)(1+sT_2) \approx 1 + sT_p$~:
\begin{equation}
T_p = T_1 + T_2 = 8~\mathrm{ms}
\qquad,\qquad
T_n = 4\,T_p = 32~\mathrm{ms}
\qquad,\qquad
T_i = 8\,K_s\,T_p^{2} = 8 \times 2{,}47 \times (0{,}008)^2 \approx 1{,}26 \times 10^{-3}~\mathrm{s}
\end{equation}

On en déduit~:
\begin{equation}
K_p = \frac{T_n}{T_i} \approx \boxed{25{,}3}
\qquad,\qquad
K_i = \frac{1}{T_i} \approx \boxed{791~\mathrm{s^{-1}}}
\end{equation}

La fonction de transfert en boucle ouverte devient alors~:
\begin{equation}
G_0(s) = G_s(s)\,G_R(s) = \frac{K_s\,(1 + s\,T_n)}{s^{2}\,T_i\,(1 + s\,T_p)}
= \frac{2{,}47\,(1 + 0{,}032\,s)}{s^{2}\cdot 1{,}26\times10^{-3}\,(1 + 0{,}008\,s)}
= \frac{1+4T_p s}{8\,T_p^{2}\,s^{2}\,(1+T_p s)}
\end{equation}
et la fonction de transfert en boucle fermée s'écrit
$G_{cf}(s) = G_0(s)/(1+G_0(s))$.

\textbf{Filtre de consigne.} La fonction $G_{cf}$ contient le zéro du PI,
$(1+sT_n)$, provenant de $G_{0}$. Ce zéro provoque un dépassement important sur
la réponse à la consigne. Pour annuler cet effet, un filtre est placé en amont
du comparateur~:
\begin{equation}
G_{lc}(s) = \frac{1}{1 + s\,T_{lc}}\ ,\qquad T_{lc} \simeq T_n = 32~\mathrm{ms}
\end{equation}
Le pôle de ce filtre annule le zéro de $G_{cf}$.

\textbf{Temps de réponse.} Pour un comportement optimal, la relation de
l'annexe~6.A donne $t_n\,\omega_c = 4{,}20$, avec $\omega_c = 1/T_p$. On en
déduit~:
\begin{equation}
t_n = 4{,}2\;T_p = 4{,}2 \times 8~\mathrm{ms} \approx \boxed{33{,}6~\mathrm{ms}}
\end{equation}

\begin{center}
\begin{tabular}{cccccc}
\toprule
$T_p$ & $T_n$ & $T_i$ & $K_p$ & $K_i$ & $T_{lc}$ \\
\midrule
$8~\mathrm{ms}$ & $32~\mathrm{ms}$ & $1{,}26~\mathrm{ms}$ & $25{,}3$ & $791~\mathrm{s^{-1}}$ & $32~\mathrm{ms}$ \\
\bottomrule
\end{tabular}
\end{center}

% -------------------------------------------------
\subsection{Question B: Régulateur PI pour 5~\% de dépassement et 0,5~s}
\label{sec:B}

Ce point n'est pas traité. Les notions nécessaires pour dimensionner un
régulateur imposant un comportement autre que le comportement « optimal » du
point~A (méthode du lieu des pôles) n'ont pas encore été vues en cours~: il est
donc normal de ne pas disposer, à ce stade, des outils pour y répondre.

Un seul régulateur est donc disponible~; la comparaison de deux régulateurs
demandée au point~C n'est par conséquent pas possible.

% -------------------------------------------------
\subsection{Question C: Vérification en simulation}

Le schéma de simulation (figure~\ref{fig:simulation}), construit sous
\textit{Simulink}, reprend la figure~\ref{fig:boucle} sans la perturbation
$K_{per}$. Il comprend un bloc \textit{Step} de $0{,}01~\mathrm{m}$
(à $t = 1~\mathrm{s}$), le gain $K_{mes}$, le filtre $G_{lc}$, un PI
($P = 25$, $I = 790$), $G_{sva}$, $G_{vey}$ et la mesure $G_{mes}$ dans la
boucle de retour. Deux scopes relèvent la position $y$ et la sortie du
régulateur.

\begin{figure}[H]
\centering
\includegraphics[width=1\textwidth]{images/schema_Simulation.png}
\caption{Schéma de simulation.}
\label{fig:simulation}
\end{figure}

\begin{figure}[H]
\centering
\includegraphics[width=1\textwidth]{images/C.png}
\caption{Réponse simulée de la position $y$ à un saut de consigne de
$0{,}01$~m avec le PI du point~A, avec (orange) et sans (bleu) filtre de
consigne.}
\label{fig:C}
\end{figure}

\begin{table}[H]
\centering
\begin{tabular}{lc}
\toprule
Grandeur & Lecture du scope \\
\midrule
Valeur finale                          & $10{,}0~\mathrm{mm}$ \\
Valeur maximale du pic                 & $\approx 10{,}85~\mathrm{mm}$ \\
\textbf{Dépassement}                   & $\mathbf{\approx 8{,}5~\%}$ \\
Temps de montée $10 \rightarrow 90~\%$ & $\approx 34~\mathrm{ms}$ \\
Temps de montée $1 \rightarrow 99~\%$  & $46{,}4~\mathrm{ms}$ \\
\textbf{Temps de réponse à $5~\%$}     & $\mathbf{\approx 86~\mathrm{ms}}$ \\
\bottomrule
\end{tabular}
\caption{Grandeurs caractéristiques de la réponse simulée (avec filtre de consigne).}
\label{tab:C}
\end{table}
Le premier constat est la différence de dépassement entre les deux réponses~: sans filtre de consigne, le dépassement atteint environ $50~\%$, alors qu'avec le filtre, il est réduit à environ $8~\%$. Le filtre de consigne est donc efficace pour réduire le dépassement.\\
Avec le filtre de consigne, la réponse (figure~\ref{fig:C}) atteint un pic
d'environ $10{,}85~\mathrm{mm}$ à environ $70~\mathrm{ms}$ après le saut, soit
un dépassement d'environ $8~\%$, puis converge sans erreur statique vers
$10~\mathrm{mm}$. Le scope indique un temps de montée de $46{,}4~\mathrm{ms}$,
les seuils étant configurés à $1~\%$ et $99~\%$ comme demandé par l'assistant de
laboratoire. Sur les seuils usuels de $10~\%$ et $90~\%$, la lecture graphique
donne environ $34~\mathrm{ms}$. La position reste dans la bande $\pm 5~\%$ à
partir d'environ $90~\mathrm{ms}$.

\textbf{Tension nécessaire en sortie de régulateur.} Pour ce saut de
$0{,}01~\mathrm{m}$, la sortie du PI atteint au maximum environ
$\mathbf{5{,}7~V}$, environ $24~\mathrm{ms}$ après le saut
(figure~\ref{fig:C2}). Cette valeur reste inférieure à la limite de
$\pm 10~\mathrm{V}$ des canaux analogiques de l'ESP32~: le régulateur peut donc
être utilisé sans saturation. Le filtre de consigne joue également un rôle vis-à-vis de la tension à l'actionneur~: sans lui, la sortie du PI vaut, dès l'instant du saut, 14V,
soit au-delà de $\pm 10~\mathrm{V}$, et le dépassement passe d'environ $8~\%$ à
environ $50~\%$.

\begin{figure}[H]
\centering
\includegraphics[width=0.92\textwidth]{images/C_filtre_tension.png}
\caption{Tension en sortie du régulateur, avec (orange) et sans (bleu) filtre
de consigne.}
\label{fig:C2}
\end{figure}

% -------------------------------------------------
\subsection{Question D: Implémentation sur l'installation réelle}

Le régulateur A ($K_p = 25{,}31$, $K_i = 791~\mathrm{s^{-1}}$, filtre de
consigne $T_{lc} = 32~\mathrm{ms}$) est implanté dans le schéma
\textit{Simulink} de l'installation (figure~\ref{fig:D1}), avec
$T_E = 1~\mathrm{ms}$, valeur faible devant $T_p = 8~\mathrm{ms}$. Le port COM
de l'ESP32 est identifié via le gestionnaire de périphériques, puis renseigné
dans les \textit{Hardware Settings}~; l'acquisition est réalisée en mode
\textit{Monitor \& Tune}. Le schéma, fourni par l'enseignant, est organisé en
zones~:

\begin{itemize}
\item \textbf{consigne} (jaune)~: profil de position obtenu en sommant des
rampes et des sauts, converti en volts ($K_{mes}$) avec un décalage
\texttt{pos\_offset}~;
\item \textbf{régulation} (orange)~: filtre d'entrée $1/(0{,}032\,s+1)$
(c'est $G_{lc}$), puis PI~; la « tension commandée » est relevée directement en
sortie du PI, \emph{avant} la conversion analogique~;
\item \textbf{installation} (cyan)~: sorties analogiques (DAC), limitées à
$\pm 10~\mathrm{V}$~;
\item \textbf{acquisition} (rose)~: entrées analogiques, puis filtre logiciel
$1/(0{,}005\,s+1)$ sur le capteur, qui reproduit $G_{mes}$~;
\item \textbf{affichage} (vert)~: retour en mètres par $1/K_{mes}$.
\end{itemize}

\begin{figure}[H]
\centering
\includegraphics[width=1\textwidth]{images/D1.png}
\caption{Schéma \textit{Simulink} implanté sur le microcontrôleur cible (ESP32)
pour le pilotage du vérin hydraulique.}
\label{fig:D1}
\end{figure}

Le profil de consigne comprend une rampe de $0$ à $0{,}05~\mathrm{m}$
($0{,}1~\mathrm{m/s}$, de $0{,}5$ à $1~\mathrm{s}$), un saut de $+0{,}05~\mathrm{m}$
à $3~\mathrm{s}$, un saut de $-0{,}05~\mathrm{m}$ à $6~\mathrm{s}$, puis une
rampe descendante vers $0$ ($0{,}025~\mathrm{m/s}$, de $8{,}5$ à
$10{,}5~\mathrm{s}$). Le résultat est présenté à la figure~\ref{fig:D2}.

\begin{figure}[H]
\centering
\includegraphics[width=0.74\textwidth]{images/D2.png}
\caption{Réponse mesurée sur l'installation réelle avec le régulateur A
($T_E = 1$~ms)~: position désirée et mesurée (haut), tension en sortie du
régulateur avant le DAC (bas).}
\label{fig:D2}
\end{figure}

Les valeurs suivantes sont relevées sur la figure~\ref{fig:D2} (précision
d'environ $1~\mathrm{mm}$ et $2~\mathrm{V}$)~:

\begin{center}
\begin{tabular}{lll}
\toprule
Phase de la consigne & Position mesurée & Sortie du régulateur \\
\midrule
Rampe $0 \rightarrow 0{,}05$~m, $0{,}1$~m/s & suivie sans écart notable & quelques volts, $< 10$~V \\
Saut $+0{,}05$~m ($t = 3$~s) & pic $0{,}135$~m ($t \approx 3{,}4$~s), soit $\approx 70~\%$ du saut, & $+175$~V puis $-111$~V \\
 & 5~\% atteint après $\approx 1{,}2$~s & \\
Saut $-0{,}05$~m ($t = 6$~s) & creux $0{,}011$~m ($t \approx 6{,}4$~s), soit $\approx 78~\%$ du saut, & $-173$~V puis $+104$~V \\
 & 5~\% atteint après $\approx 1{,}1$~s & \\
Rampe $0{,}05 \rightarrow 0$~m, $0{,}025$~m/s & suivie sans écart notable & $\ll 10$~V \\
\bottomrule
\end{tabular}
\end{center}

\textbf{Effet de la pente de consigne.} Lorsque la consigne présente une pente
faible, la position suit correctement la consigne. Lorsqu'elle présente une
pente raide (sauts de $0{,}05~\mathrm{m}$), la sortie du régulateur dépasse
largement la limite de tension de l'actionneur ($\pm 10~\mathrm{V}$) et la commande sature. L'intégrateur s'emballe~: il continue de se charger tant que la position n'a pas atteint la consigne, puis doit se décharger. Il en résulte un dépassement de $70~\%$ et des oscillations sur environ une seconde, avant que le système ne se stabilise. La limite est donc celle de l'actionneur et non celle du régulateur.

\textbf{Remarque.} Le régulateur est dimensionné avec un modèle linéaire qui
suppose l'actionneur non saturé. Ce dimensionnement n'est donc valable que pour
des consignes assez douces~: en pratique, on limitera la pente ou l'amplitude
de la consigne.

% -------------------------------------------------
\subsection{Question E: Un régulateur PID ou P serait-il envisageable~?}

\textbf{Régulateur PID.} Avec le PI, le système sature déjà. Ajouter une action
dérivée ne ferait qu'augmenter la vitesse de la commande, et donc la
saturation~: un régulateur plus rapide réclamerait des commandes encore plus
grandes, on saturerait davantage et plus longtemps, et l'intégrateur
s'emballerait encore plus. Le PID n'est donc pas une solution envisageable.

\textbf{Régulateur P.} Le système à régler est de type~1~: un régulateur
proportionnel suffit à annuler l'erreur statique de position pour un saut de
consigne. En revanche, comme des perturbations
de débit agissent sur le système, l'erreur statique ne serait pas nulle et le P
ne peut pas la corriger~: il subsisterait toujours un écart entre la position
désirée et la position mesurée. Le P n'est donc pas une solution envisageable.

% =================================================
\section{Conclusion}

Le système du vérin hydraulique étant de type~1, le PI a été dimensionné par le
réglage en correspondance~: $T_p = 8~\mathrm{ms}$,
$T_n = 4T_p = 32~\mathrm{ms}$, $T_i = 8K_sT_p^2 \approx 1{,}26~\mathrm{ms}$, soit
$K_p \approx 25{,}3$ et $K_i \approx 791~\mathrm{s^{-1}}$, avec un filtre de
consigne $T_{lc} \simeq T_n$ qui annule l'effet du zéro du PI. Le temps
caractéristique estimé, $t_n = 4{,}2\,T_p = 33{,}6~\mathrm{ms}$, est confirmé en
simulation (temps de montée d'environ $34~\mathrm{ms}$, dépassement d'environ
$8~\%$, tension maximale de $5{,}7~\mathrm{V}$ pour un saut de
$0{,}01~\mathrm{m}$). Le dimensionnement pour un comportement imposé (point~B)
n'a pas pu être traité, la matière n'ayant pas encore été vue en cours.

Sur l'installation réelle, le régulateur suit correctement les rampes lentes,
mais la commande sature dès que la consigne présente une pente raide (saut de
$0{,}05~\mathrm{m}$)~: l'intégrateur s'emballe, ce qui donne un dépassement de
$70~\%$ et des oscillations sur environ une seconde. La limite est celle de
l'actionneur, ce qui exclut l'usage d'un PID, plus rapide, et incite plutôt à
limiter la pente de la consigne ou à prévoir une protection anti-windup.

\end{document}